\scanpage{12}
\mspage{5}
% ll. 1–3
$A\in\operatorname{Ob}\mathrm{Cat}$ (je pense au cas $A=\widetilde\Delta$)

$\pi\in\operatorname{Ob}A^\wedge$, $Z=(A_{/\pi})^{\mathrm o}\simeq\pi\backslash\underbrace{A^{\mathrm o}}_B=\pi\backslash B$

$[\pi:\underbrace{A^{\mathrm o}}_B\longrightarrow\mathrm{Ens}]$ On s'intéresse à $Z^\wedge$, et à la sous-catégorie pleine de $Z^\wedge$ \deleted{\unsure{engendré|engendr.|engendrée}} stable par $\varprojlim$ finies, engendrée par $Z$. J'ai en vue un candidat : pour tout $b\in\operatorname{Ob}B$, considérons $Z_b\simeq\pi(\uline{b})$, et pour tte partie \inserted{\uline{non vide}} $E$ de $\pi(b)$
% ll. 8–12
\diagram{\begin{tikzcd}[ampersand replacement=\&,row sep=small]
 Z\arrow[d]\& Z_b\simeq\pi(b)\\ B\ni b\&
\end{tikzcd}}
\[
 E\in\mathcal P^*(\pi(b))
\]
considérons l'élément
\[
 \widetilde E=\coprod_{a\in E}a\qquad(\text{somme prise \unsure{ds|dans} }Z^\wedge)
\]
% ll. 13–16
Si $u:b\longrightarrow b'$ ds $B$, on a \inserted{$u_*=\pi(u):$} $\underset{\substack{\rotatebox{90}{$=$}\\\pi(b)}}{Z_b}\longrightarrow\underset{\substack{\rotatebox{90}{$=$}\\\pi(b')}}{Z_{b'}}$ et si $E\subset Z_b$, $E'\subset Z_{b'}$ sont tels que $u_*(E)\subset E'$, alors $u$ induit
\[
 u_*^{E,E'}:\widetilde E\longrightarrow\widetilde E'
\]
% ll. 17–19
Ainsi, le foncteur composé
\[
 \widetilde\pi:b\longmapsto\mathcal P^*(\pi(b)):\quad B\xrightarrow{\pi}\mathrm{Ens}\xrightarrow{\mathcal P^*}\mathrm{Ord}
\]
définit une catégorie $\widetilde\pi\backslash B$, \deleted{\unsure{ayant d|ayant q|dont}} $=\unsure{S|s}$
% archive sheet number: 11
